\documentclass[11pt]{article}
\usepackage[T1]{fontenc}
\usepackage{lmodern}
\usepackage{amsmath,amssymb,amsthm,mathtools}
\usepackage[margin=1in]{geometry}
\usepackage{microtype}
\usepackage[colorlinks=true,linkcolor=blue,citecolor=blue,urlcolor=blue]{hyperref}
\newtheorem{theorem}{Theorem}[section]
\newtheorem{lemma}[theorem]{Lemma}
\newtheorem{proposition}[theorem]{Proposition}
\newtheorem{remark}[theorem]{Remark}
\DeclareMathOperator{\Tr}{Tr}
\DeclareMathOperator{\tr}{tr}
\DeclareMathOperator{\Spec}{Spec}
\DeclareMathOperator{\Var}{Var}
\newcommand{\E}{\mathbb E}
\newcommand{\Prob}{\mathbb P}
\newcommand{\C}{\mathbb C}
\newcommand{\M}{\mathrm M}
\newcommand{\HS}{\mathrm{HS}}
\newcommand{\id}{\mathrm{id}}
\newcommand{\norm}[1]{\left\|#1\right\|}
\title{Strong convergence for a repeated GUE interaction\\ on three tensor legs}
\author{Kaiwen Jin\\ \small\href{mailto:yui.ui.siki@gmail.com}{yui.ui.siki@gmail.com}}
\date{13 October 2026}
\hypersetup{pdftitle={Strong convergence for a repeated GUE interaction on three tensor legs},
pdfauthor={Kaiwen Jin}}
\begin{document}
\maketitle

\begin{abstract}
Let $X_N$ be a normalized Gaussian unitary ensemble matrix of size $N^2$.
We place the same realization on two adjacent pairs of tensor factors,
as $X_N\otimes I_N$ and $I_N\otimes X_N$ on $(\C^N)^{\otimes3}$.
For every fixed complex noncommutative polynomial, its operator norm
converges in probability to the norm of that polynomial in two free
variance-one semicircular elements. The argument computes the covariance
of the dependent pair and applies existing Gaussian spectral comparison
and operator-valued free convergence results. A direct Gaussian variance
estimate supplies the trace convergence needed for the norm lower bound.
\end{abstract}

\section{Introduction}

A local interaction may be used repeatedly in a quantum spin system.
The resulting operators act on different tensor factors, while their
random coefficients come from one matrix. Collins and Yuan
\cite[Section~1.1 and Section~5(2)]{CY}
describe this repeated-interaction question and study tensor embeddings
with independent local GUE matrices. The more general tensor models in
Chen, Garza-Vargas and van Handel \cite[Section~9.4, Theorem~9.8]{CGV}
also use independent local matrices. Here we retain the common Gaussian
coefficients and consider two overlapping ordered embeddings.

For each positive integer $N$, let $X_N$ be a normalized Hermitian GUE
matrix on $\C^N\otimes\C^N$. Thus its diagonal entries have variance
$N^{-2}$, and its upper-triangular complex entries satisfy
$\E|(X_N)_{ab}|^2=N^{-2}$. Define
\begin{equation}\label{eq:model}
 A_N=X_N\otimes I_N,\qquad B_N=I_N\otimes X_N
 \quad\text{on }(\C^N)^{\otimes3}.
\end{equation}
Both expressions use exactly the same $X_N$, with the displayed tensor
factor order. Let $s_1,s_2$ be centered free semicircular elements of
variance one in their reduced free-product $C^*$-algebra, equipped with
its faithful tracial representation.

\begin{theorem}\label{thm:main}
For every fixed polynomial $P\in\C\langle x,y\rangle$,
\[
 \norm{P(A_N,B_N)}
 \xrightarrow[N\to\infty]{\Prob}
 \norm{P(s_1,s_2)}.
\]
All norms are operator norms, and the norm on the right is the norm in
the reduced free-product $C^*$-algebra.
\end{theorem}

This proves the three-site, two-interaction case of the repeated-matrix
question. The conclusion includes non-self-adjoint polynomials and
does not impose a degree restriction on a fixed polynomial. Larger
families of repeated interactions are outside the present result.

The general comparison and convergence results are due to
Bandeira, Boedihardjo and van Handel \cite{BBvH}, Gao and
Kunnawalkam Elayavalli \cite{GK}, and Jekel, Lee, Nelson and Pi \cite{JLNP}.
Our application consists of the covariance calculation for the shared
realization and the passage through complete linear matrix pencils.
The joint entry covariance has norm $2/N$, while every fixed covariance
expression formed from the identity tends to its scalar free value.
These facts permit comparison with the free pair. We keep the constant
matrix coefficient inside the Gaussian comparison, and use the
deterministic linearization of Haagerup and Thorbj{\o}rnsen \cite{HT}.
Expected trace convergence and a variance bound then give the reverse
norm inequality.

We write $\Tr$ for unnormalized trace and $\tr_k=k^{-1}\Tr$.
Hilbert--Schmidt norms are unnormalized:
$\norm{R}_{\HS}^2=\Tr(R^*R)$. The proof estimates are made for $N\ge2$;
a finite initial dimension does not affect the limit.

\section{Covariance of the shared realization}

Put $m=N^2$ and $n=N^3$. Choose a Hermitian orthonormal basis
$(H_\alpha)_{\alpha=1}^{N^4}$ of $\M_m(\C)$ for
$\Tr(H_\alpha H_\beta)=\delta_{\alpha\beta}$.
For the given matrix, the coordinates
$g_\alpha=N\Tr(H_\alpha X_N)$ are independent standard real Gaussians.
Hence
\begin{equation}\label{eq:coefficients}
 X_N=\frac1N\sum_\alpha g_\alpha H_\alpha,\qquad
 C_{1\alpha}=\frac{H_\alpha\otimes I_N}{N},\qquad
 C_{2\alpha}=\frac{I_N\otimes H_\alpha}{N}.
\end{equation}
With $Y_{1,N}=A_N$ and $Y_{2,N}=B_N$, both coordinates are given by
\begin{equation}\label{eq:joint-gaussian}
 Y_{i,N}=\sum_\alpha g_\alpha C_{i\alpha}.
\end{equation}
In particular, the scalar index is shared in the two sums.

The covariance maps on $\M_n(\C)$ are
\[
 \eta_{ij,N}(R)=\E[Y_{i,N}RY_{j,N}]
             =\sum_\alpha C_{i\alpha}RC_{j\alpha}.
\]
The block map $R\mapsto(\eta_{ij,N}(R))_{i,j=1}^2$ is completely
positive: stack $C_{1\alpha}$ and $C_{2\alpha}$ to form its Kraus
operators. Since these matrices are Hermitian, cyclicity of trace gives
\begin{equation}\label{eq:trace-symmetry}
 \tr_n(\eta_{ij,N}(R)T)=\tr_n(R\eta_{ji,N}(T)).
\end{equation}
All component maps are normal in finite dimension.

\begin{lemma}\label{lem:covariance}
For $i,j\in\{1,2\}$,
\begin{align}
 \sum_\alpha C_{i\alpha}^2&=I_n,
 &\norm{\eta_{ij,N}(R)}&\le\norm R, \label{eq:map-bounds}\\
 \eta_{11,N}(I_n)&=\eta_{22,N}(I_n)=I_n,
 &\eta_{12,N}(I_n)&=N^{-2}U_N,\quad
 \eta_{21,N}(I_n)=N^{-2}U_N^*, \label{eq:identity-values}\\
 \Tr(C_{i\alpha}C_{i\beta})&=N^{-1}\delta_{\alpha\beta}.
 &&\label{eq:gram}
\end{align}
Here $U_N$ is the cyclic permutation unitary
$U_N(e_a\otimes e_b\otimes e_c)=e_b\otimes e_c\otimes e_a$.
The joint entry-covariance operator, in the unnormalized
Hilbert--Schmidt convention, has norm $2/N$.
\end{lemma}

\begin{proof}
Basis completeness and Hermitian symmetry give
\[
 \sum_\alpha (H_\alpha)_{ab}(H_\alpha)_{cd}
       =\delta_{ad}\delta_{bc},\qquad
 \sum_\alpha H_\alpha^2=mI_m.
\]
The second identity proves the first part of \eqref{eq:map-bounds}.
The coefficient row $(C_{i\alpha})_\alpha$ has norm one, as does
the corresponding adjoint column. Factoring the map as row, diagonal
copies of $R$, and column proves the map bound.

The entry covariance of the original matrix is
$\E[(X_N)_{ab}(X_N)_{cd}]=N^{-2}\delta_{ad}\delta_{bc}$.
Using the ordered embeddings, the row $(i,j,k)$ and column $(l,p,q)$
entry of $\E[A_NB_N]$ is
\[
 N^{-2}\delta_{ip}\delta_{jq}\delta_{kl}.
\]
This is $N^{-2}U_N$; taking adjoints gives the other cross term.
The diagonal identity values follow from the first part of the lemma.

Equation \eqref{eq:gram} follows from
$N^{-2}\Tr(H_\alpha H_\beta)\Tr(I_N)=N^{-1}\delta_{\alpha\beta}$.

Let $V_N$ have columns
$(\operatorname{vec}C_{1\alpha},\operatorname{vec}C_{2\alpha})$
in the direct sum of two Hilbert--Schmidt spaces.
The joint covariance operator is $T_N=V_NV_N^*$.
Thus
\begin{equation}\label{eq:choi}
 V_N^*V_N=\frac2N I_{N^4},\qquad \norm{T_N}=\frac2N.
\end{equation}
This is also the Choi operator in the convention of
\cite[Definition~1.5 and Remark~1.6]{JLNP}.
\end{proof}

For the limiting scalar base algebra $\C$, put
$\eta_{ij}(z)=\delta_{ij}z$.
A deterministic covariance expression means a fixed finite expression
formed from scalars and the identity using sums, products, adjoints
and the maps $\eta_{ij,N}$. Its coefficients and expression do not
depend on $N$.

\begin{lemma}\label{lem:base-convergence}
For each such expression $f$, there is a constant $K_f$ such that
\begin{equation}\label{eq:scalar-approximation}
 \norm{f(\eta_N,I_n)-f(\eta,1)I_n}\le K_fN^{-2},
 \qquad N\ge2.
\end{equation}
Consequently both its normalized trace and its norm converge to their
scalar values.
\end{lemma}

\begin{proof}
Induct over the finite expression. Scalars and the identity give the
base cases. Sums and adjoints preserve the estimate.
For products use
$R_NS_N-cdI=(R_N-cI)S_N+c(S_N-dI)$;
the preceding estimates bound $S_N$ uniformly.
For a map application, write $R_N=cI_n+E_N$. Lemma~\ref{lem:covariance}
gives
\[
 \norm{\eta_{ij,N}(R_N)-\delta_{ij}cI_n}
 \le\norm{E_N}+|c|N^{-2}\mathbf1_{\{i\ne j\}}.
\]
This proves the induction. The trace conclusion follows from
$|\tr_n(E_N)|\le\norm{E_N}$.
\end{proof}

The cross map in Lemma~\ref{lem:covariance} is only small at the
identity. The argument uses its uniform norm bound for general inputs;
it does not require that its full operator norm tend to zero.

\section{The spectral upper bound}

Let $(h_\alpha)$ be a free variance-one semicircular family in a
faithful tracial von Neumann algebra, and define
\begin{equation}\label{eq:free-replacement}
 S_{i,N}=\sum_\alpha C_{i\alpha}\otimes h_\alpha.
\end{equation}
The same $h_\alpha$ occurs in both components.
By \cite[Proposition~1.8]{JLNP}, this is an $\M_n$-valued
semicircular family with covariance $\eta_N$.

\begin{lemma}\label{lem:free-convergence}
For every fixed complex polynomial $P$,
\[
 \norm{P(S_{1,N},S_{2,N})}\longrightarrow\norm{P(s_1,s_2)}.
\]
\end{lemma}

\begin{proof}
Use \cite[Corollary~1.4 and Examples~2.5(3)]{GK}, in the tracial
form explained in \cite[Theorem~4.3]{JLNP}.
The finite base algebras are $\M_n$, the base tuples are $I_n$,
and the limiting base is $\C$ generated by $1$.
The covariance is block completely positive and trace-symmetric
by \eqref{eq:trace-symmetry}; its diagonal value at $1$ is bounded by one.
Lemma~\ref{lem:base-convergence} proves strong convergence of all
deterministic covariance expressions, including normalized traces.

The norm tests in \cite[Examples~2.5(3)]{GK} can also be read directly:
for fixed $p,q$ in the base variable,
\[
 \norm{p(I_n)-\eta_{ij,N}(q(I_n))}
 \longrightarrow |p(1)-\delta_{ij}q(1)|.
\]
The generation condition is imposed on the limiting base, not on
the finite algebras. The faithful product trace on
$\M_n\otimes W^*(h_\alpha)$ provides the tracial realization of
\eqref{eq:free-replacement}. The chosen limit is the target free pair,
with scalar covariance $\delta_{ij}$.
\end{proof}

We need this convergence for matrix pencils as well as scalar
polynomials. The following argument also covers continuous spectra.

\begin{lemma}\label{lem:free-spectra}
Fix $d\ge1$ and self-adjoint $a_0,a_1,a_2\in\M_d(\C)$, and set
\[
 Q_{N,\mathrm{free}}=a_0\otimes1+
                   a_1\otimes S_{1,N}+a_2\otimes S_{2,N},\qquad
 Q=a_0\otimes1+a_1\otimes s_1+a_2\otimes s_2.
\]
There is a sequence $\delta_N\to0$ such that
\[
 \Spec(Q_{N,\mathrm{free}})
       \subseteq\Spec(Q)+[-\delta_N,\delta_N].
\]
\end{lemma}

\begin{proof}
Put $\mathcal A_N=C^*(1,S_{1,N},S_{2,N})$ and
$\mathcal A=C^*(1,s_1,s_2)$.
The norm in $\prod_N\mathcal A_N/\bigoplus_N\mathcal A_N$ is the norm
limsup, where the direct sum denotes norm-null sequences.
Indeed norm-null corrections preserve the limsup, and cancelling
finite initial segments gives the reverse inequality.
Lemma~\ref{lem:free-convergence} therefore makes
\[
 P(s_1,s_2)\longmapsto[(P(S_{1,N},S_{2,N}))_N]
\]
a well-defined isometry on the polynomial algebra. It extends to a
unital $*$-embedding of $\mathcal A$ into this quotient.
Apply the map entrywise in $\M_d$. Matrix boundedness and
norm-nullness are equivalent to the corresponding properties of
their finitely many entries.

If the claimed spectral containment failed, uniform boundedness would
give a subsequence of spectral points $\lambda_N$ tending to a real
$\lambda\notin\Spec(Q)$. The inverse of $\lambda-Q$ maps to an inverse
of the matrix sequence class. Choose a bounded representative $R_N$.
Both inverse residuals tend to zero. For all sufficiently large $N$,
their norms are less than $1/2$, so the Neumann series gives an actual
two-sided inverse of $\lambda-Q_{N,\mathrm{free}}$, with norm at most
$2\sup_N\norm{R_N}$.

Inversion sends $\lambda_N$ to the spectral point
$(\lambda-\lambda_N)^{-1}$ of that inverse. Its norm must tend to
infinity, a contradiction. If $\lambda=\lambda_N$, invertibility
already fails. This uses spectral points, without an eigenspace
assumption. The sequence-quotient argument is the deterministic
mechanism in \cite[pp.~758--760]{HT}.
\end{proof}

Now consider the Gaussian pencil
\begin{equation}\label{eq:full-pencil}
 Q_N=a_0\otimes I_n+a_1\otimes A_N+a_2\otimes B_N
     =a_0\otimes I_n+\sum_\alpha g_\alpha L_{\alpha,N},
 \quad
 L_{\alpha,N}=a_1\otimes C_{1\alpha}+a_2\otimes C_{2\alpha}.
\end{equation}
Its free replacement is precisely $Q_{N,\mathrm{free}}$.
For a self-adjoint Gaussian pencil
$Y=L_0+\sum_\alpha g_\alpha L_\alpha$ of size $D$,
the parameters in \cite[Section~2.1]{BBvH} are
\[
 \sigma(Y)^2=\norm{\sum_\alpha L_\alpha^2},\qquad
 v(Y)^2=\sup_{\Tr(R^*R)\le1}\sum_\alpha|\Tr(L_\alpha R)|^2.
\]
The source's parameter $\sigma_*$ satisfies $\sigma_*\le v$.
Consequently \cite[Theorem~2.1]{BBvH} gives, for all $t\ge0$,
\begin{equation}\label{eq:bbvh}
 \Prob\!\left\{\Spec(Y)\subseteq\Spec(Y_{\mathrm{free}})
  +C\big[v(Y)^{1/2}\sigma(Y)^{1/2}(\log D)^{3/4}
                +v(Y)t\big][-1,1]\right\}\ge1-e^{-t^2}.
\end{equation}
Here $Y_{\mathrm{free}}=L_0\otimes1+\sum_\alpha L_\alpha\otimes h_\alpha$.
The deterministic shift is unchanged, and is not included in either
noise parameter.

For \eqref{eq:full-pencil}, let
$M_a=\norm{a_1}+\norm{a_2}$ and
$K_a=2(\norm{a_1}_{\HS}^2+\norm{a_2}_{\HS}^2)$.
Lemma~\ref{lem:covariance} gives
\[
 \sum_\alpha L_{\alpha,N}^2
 =a_1^2\otimes I_n+a_2^2\otimes I_n+
 N^{-2}(a_1a_2\otimes U_N+a_2a_1\otimes U_N^*),
\]
so $\sigma(Q_N)\le M_a$.
Bessel's inequality using \eqref{eq:gram} gives, for each component,
\[
 \sum_\alpha|\Tr((a_i\otimes C_{i\alpha})R)|^2
 \le N^{-1}\norm{a_i}_{\HS}^2\norm{R}_{\HS}^2.
\]
Applying two-term Cauchy--Schwarz to the shared coefficient sum yields
$v(Q_N)^2\le K_a/N$.

Choose $t_N=(\log(dN^3))^{3/2}$ in \eqref{eq:bbvh}.
Then, with probability at least $1-e^{-(\log(dN^3))^3}$,
\begin{equation}\label{eq:gaussian-free-inclusion}
 \Spec(Q_N)\subseteq\Spec(Q_{N,\mathrm{free}})+[-r_N,r_N],
\end{equation}
where
\[
 r_N\le C\left[
 K_a^{1/4}M_a^{1/2}N^{-1/4}(\log(dN^3))^{3/4}
 +K_a^{1/2}N^{-1/2}(\log(dN^3))^{3/2}\right]\longrightarrow0.
\]
The failure probabilities are summable.
Borel--Cantelli and Lemma~\ref{lem:free-spectra} give, for each fixed
pencil and each $\varepsilon>0$,
\begin{equation}\label{eq:final-spectral-inclusion}
 \Spec(Q_N)\subseteq\Spec(Q)+[-\varepsilon,\varepsilon]
 \quad\text{eventually, almost surely}.
\end{equation}
The comparison was applied to the whole pencil, including $a_0$;
no addition of separate spectral containments is used.

\begin{proposition}\label{prop:upper}
On a probability space carrying the GUE sequence,
\[
 \limsup_{N\to\infty}\norm{P(A_N,B_N)}
       \le\norm{P(s_1,s_2)}
 \quad\text{almost surely for every }P\in\C\langle x,y\rangle.
\]
\end{proposition}

\begin{proof}
Intersect the events in \eqref{eq:final-spectral-inclusion} for all
positive integer $d$ and all self-adjoint coefficient matrices with
entries in $\mathbb Q+i\mathbb Q$. This is a countable intersection.
Coordinate pencils give almost-sure boundedness of $A_N,B_N$.

Fix such an outcome and form
$\mathcal C=\prod_N\M_{N^3}/\bigoplus_N\M_{N^3}$,
with coordinate classes $y_1,y_2$.
For any rational self-adjoint pencil, a point outside $\Spec(Q)$ is
at a fixed positive distance from every sufficiently late finite
Gaussian spectrum. The late resolvents are uniformly bounded; setting
the finitely many early resolvents to zero gives an inverse in the
matrix quotient. Thus its spectrum is contained in $\Spec(Q)$.
The deterministic \cite[Theorem~2.4]{HT} gives a unital
$*$-homomorphism from $C^*(1,s_1,s_2)$ onto $C^*(1,y_1,y_2)$,
sending $s_i$ to $y_i$.
Its contractivity and the quotient norm-limsup formula give the
conclusion for all complex polynomials.
This is also the criterion stated in \cite[Theorem~7.7]{BBvH}.
\end{proof}

\section{Trace convergence and the norm lower bound}

The expected trace theorem \cite[Lemma~3.9]{JLNP} applies to our
joint Gaussian tuple, with deterministic base data $I_n$.
Lemma~\ref{lem:covariance} supplies trace symmetry and bounded
component maps; \eqref{eq:choi} gives Choi norms tending to zero
also for every coordinate restriction.
Lemma~\ref{lem:base-convergence} gives every deterministic
covariance-expression trace limit. The identity base is independent
of the Gaussian tuple and has uniformly bounded moments.
Hence, for each word $w$,
\begin{equation}\label{eq:expected-moments}
 \E\tr_n w(A_N,B_N)\longrightarrow\tau w(s_1,s_2).
\end{equation}
Here $\tau$ is the faithful trace of the limiting free pair.

We first bound the operator-norm moments needed for a direct
Poincar\'e estimate. In \eqref{eq:bbvh}, use the coordinate pencil
$A_N$. Its parameters are $\sigma(A_N)=1$ and
$v(A_N)=N^{-1/2}$ by \eqref{eq:map-bounds} and
\eqref{eq:gram}. The free coordinate norms are uniformly bounded by
Lemma~\ref{lem:free-convergence}, and
$N^{-1/4}(\log(N^3))^{3/4}$ is bounded.
Since $\norm{A_N}=\norm{B_N}=\norm{X_N}$, constants $M,C>0$ independent
of $N\ge2$ satisfy
\begin{equation}\label{eq:coordinate-tail}
 \Prob\{\norm{X_N}>M+Ct\}\le e^{-t^2},\qquad t\ge0.
\end{equation}
Partition into $[0,M]$ and $(M+Ck,M+C(k+1)]$, $k\ge0$.
For each fixed positive integer $r$,
\[
 \sup_N\E\norm{X_N}^r
 \le M^r+\sum_{k\ge0}(M+C(k+1))^r e^{-k^2}<\infty.
\]

\begin{lemma}\label{lem:variance}
For every word $w$ of length $\ell\ge1$ there is a constant $C_\ell$
independent of $N$ such that
\[
 \Var(\tr_n w(A_N,B_N))\le C_\ell N^{-4}.
\]
Consequently the normalized traces of all words converge almost surely
to their traces in the free pair.
\end{lemma}

\begin{proof}
Write $F_N=\tr_n(Y_{i_1,N}\cdots Y_{i_\ell,N})$.
Differentiate with respect to the actual shared scalar $g_\alpha$
in \eqref{eq:joint-gaussian}. Cyclicity gives
\[
 \partial_\alpha F_N=\frac1n\sum_{t=1}^{\ell}
                  \Tr(C_{i_t,\alpha}W_t),
\]
where $W_t$ is the cyclic product of the remaining $\ell-1$ factors.
Using Cauchy--Schwarz over the positions and Bessel's inequality
for \eqref{eq:gram},
\[
 \sum_\alpha|\partial_\alpha F_N|^2
 \le\frac{\ell}{Nn^2}\sum_t\norm{W_t}_{\HS}^2
 \le\frac{\ell^2}{Nn}\norm{X_N}^{2\ell-2}.
\]
The Gaussian Poincar\'e inequality
\cite[Proposition~4.1 and Corollary~4.2]{HT}
applies to this complex polynomial and its square-integrable gradient.
The uniform moments just proved and $n=N^3$ give the variance bound.
Constant words have zero variance.

For each positive $\varepsilon$, the squared-deviation inequality gives
\[
 \Prob\{|F_N-\E F_N|>\varepsilon\}
          \le C_\ell\varepsilon^{-2}N^{-4}.
\]
This is summable, so Borel--Cantelli gives
$F_N-\E F_N\to0$ almost surely.
Combine with \eqref{eq:expected-moments} and intersect the countably
many word events. Finite linear combinations then cover all
complex polynomial traces. No independence between different $N$
is used.
\end{proof}

\begin{proof}[Proof of Theorem~\ref{thm:main}]
Place the GUE matrices on a common probability space.
For a fixed $P$, put
$Z_N=P(A_N,B_N)^*P(A_N,B_N)$ and
$Z=P(s_1,s_2)^*P(s_1,s_2)$.
For every fixed integer $q\ge1$, Lemma~\ref{lem:variance} gives
$\tr_n(Z_N^q)\to\tau(Z^q)$ almost surely. Since $Z_N$ is positive,
\[
 \liminf_N\norm{P(A_N,B_N)}
       \ge \tau(Z^q)^{1/(2q)}.
\]
The limiting trace is faithful, so
$\tau(Z^q)^{1/q}\to\norm Z$.
Indeed the upper bound follows from $Z\le\norm Z\,1$, and a nonzero
spectral neighborhood below $\norm Z$ has positive trace, giving the
reverse bound after taking roots. The case $Z=0$ is immediate.
This is the trace-power argument of \cite[Lemma~4.1]{JLNP}.
Taking $q\to\infty$ after the fixed-$q$ limit proves
\[
 \liminf_N\norm{P(A_N,B_N)}\ge\norm{P(s_1,s_2)}.
\]
Combine with Proposition~\ref{prop:upper} to obtain the norm limit
almost surely on this coupling. It therefore holds in probability
for the prescribed GUE marginals. In the displayed model, $I_N$
is the $N$ by $N$ identity, both embeddings keep the ordered legs,
and $P$ has fixed complex coefficients independent of $N$.
\end{proof}

\begin{thebibliography}{9}
\bibitem{BBvH}
Afonso S. Bandeira, March T. Boedihardjo and Ramon van Handel,
\emph{Matrix concentration inequalities and free probability},
Invent. Math. \textbf{234} (2023), 419--487.
Version used: \href{https://arxiv.org/abs/2108.06312v2}{arXiv:2108.06312v2}.

\bibitem{CGV}
Chi-Fang Chen, Jorge Garza-Vargas and Ramon van Handel,
\emph{A new approach to strong convergence II: The classical ensembles},
\href{https://arxiv.org/abs/2412.00593v3}{arXiv:2412.00593v3} (2026).

\bibitem{CY}
Beno\^{i}t Collins and Wangjun Yuan,
\emph{Strong convergence for tensor GUE random matrices},
\href{https://arxiv.org/abs/2407.09065v1}{arXiv:2407.09065v1} (2024).

\bibitem{GK}
David Gao and Srivatsav Kunnawalkam Elayavalli,
\emph{Toeplitz exactness for strong convergence},
\href{https://arxiv.org/abs/2604.14106v3}{arXiv:2604.14106v3} (2026).

\bibitem{HT}
Uffe Haagerup and Steen Thorbj{\o}rnsen,
\emph{A new application of random matrices:
$\operatorname{Ext}(C^*_{\mathrm{red}}(\mathbb F_2))$ is not a group},
Ann. of Math. \textbf{162} (2005), 711--775.
\href{https://doi.org/10.4007/annals.2005.162.711}{doi:10.4007/annals.2005.162.711}.

\bibitem{JLNP}
David Jekel, Yoonkyeong Lee, Brent Nelson and Jennifer Pi,
\emph{Strong convergence to operator-valued semicirculars},
\href{https://arxiv.org/abs/2506.19940v3}{arXiv:2506.19940v3} (2026).
\end{thebibliography}

\section*{AI use}
OpenAI Codex assisted with deriving and checking the argument,
retrieving references, and preparing the manuscript.
The mathematical claims are submitted for independent review.

\end{document}
